\section{Numerische Umsetzung}
\label{sec:numerik}

% VERANTWORTLICH: ...

\subsection{Programmstruktur}

% Kurze Uebersicht der Dateien und wer was aufruft. Ein Absatz plus Liste
% reicht, der Code selbst gehoert nicht ins Protokoll.
% Erwaehnen: Rechnen und Plotten sind getrennt, Ergebnisse werden in
% .mat-Dateien zwischengespeichert.

\subsection{Wahl des Loesungsverfahrens}
\label{sec:solver}

% Das ist der Abschnitt, in dem die Vorlesung am direktesten geprueft wird.
% Zu beantworten:
%   - Ist das Problem steif? Das System ist skalar, ein Steifheitsverhaeltnis
%     ist damit gar nicht definiert. Daraus folgt die Wahl eines expliziten
%     Verfahrens.
%   - Warum ode45 (Dormand-Prince, Ordnung 5 mit eingebetteter Ordnung 4
%     zur Schrittweitensteuerung) und nicht ode15s.
%   - Welche Toleranzen wurden gesetzt und warum nicht die Defaults.
%   - Nachweis der Loesungsunabhaengigkeit: Rechnung mit strengerer Toleranz
%     wiederholen und zeigen, dass sich das Ergebnis nicht mehr aendert.

\subsection{Behandlung der Wetterdaten}

% Warum die Messreihen interpoliert werden muessen: der Solver waehlt seine
% Zeitschritte adaptiv und fragt Zeitpunkte zwischen den Messwerten ab.
% Warum lineare und nicht kubische Interpolation.
